% !TEX program = xelatex
\documentclass[a4paper,11pt,fontset=none]{ctexart}

\newcommand{\reporttitle}{低位蓄势板块与主升浪启动信号扫描}
\newcommand{\reportsubtitle}{7月资金高低切换 | 主力悄悄吸筹 | 五大板块即将接力}
\newcommand{\reportkicker}{ASTOCK RESEARCH NOTE}
\newcommand{\reportscope}{CHINA A-SHARES | SECTOR ROTATION}
\newcommand{\reportdate}{2026-07-07}
\newcommand{\reportdatacutoff}{2026-07-07 (Fund flow 07/06; TrendForce; broker reports; chiplist data 06/30)}
\newcommand{\reporttype}{Sector rotation deep scan}
\newcommand{\reportauthor}{AStock Agent Team}
\newcommand{\reporthouseview}{%
7月A股呈现极致的``高位兑现+低位承接''格局：上半年翻倍的AI算力/光模块遭机构集中抛售，资金系统性迁徙至低位板块。龙虎榜、北向资金、筹码数据共同指向五大低位蓄势方向：(1) 创新药——已走出独立主升，连续3日资金净流入超10亿，估值处近5年20\%分位；(2) 半导体设备/材料——北向7月净流入超110亿，订单锁至2027，全年横盘后即将启动；(3) 军工——估值近10年18\%低位，社保/北向逆势抄底，十五五刚需+军贸爆发；(4) 人形机器人零部件——超跌小票批量涨停，宇树IPO+Optimus量产双催化；(5) 氟化工/电子特气——产品涨价17--232\%，AI绑定+配额刚性，周期彻底反转。高位科技兑现资金的迁徙周期通常持续15--20个交易日，当前处于早期阶段。%
}
\newcommand{\reportquality}{Fund flow data (07/06-07/07 confirmed); Northbound flow (confirmed); chiplist concentration (confirmed 06/30); broker reports (06-07/2026). High confidence / multi-source.}
\newcommand{\reportdisclaimer}{本报告不构成任何投资建议。}
% Reusable IB-style preamble for XeLaTeX equity research reports
% Usage: % Reusable IB-style preamble for XeLaTeX equity research reports
% Usage: % Reusable IB-style preamble for XeLaTeX equity research reports
% Usage: \input{preamble} after \documentclass[a4paper,11pt,openany,fontset=none]{ctexrep}

% ============ 页面布局 ============
\usepackage[top=2cm, bottom=2cm, left=2cm, right=2cm]{geometry}
\setlength{\parskip}{0.3em}
\setlength{\parindent}{2em}
\renewcommand{\baselinestretch}{1.15}

% ============ 字体（macOS） ============
\setCJKmainfont[BoldFont={Heiti SC}]{STSong}
\setCJKsansfont{Heiti SC}
\setCJKmonofont{Heiti SC}
\setmainfont{Times New Roman}

% ============ 颜色（IB Navy/Grey palette） ============
\usepackage{xcolor}
\definecolor{navy}{HTML}{003366}
\definecolor{steelgrey}{HTML}{4A5568}
\definecolor{lightgrey}{HTML}{F7FAFC}
\definecolor{riskred}{HTML}{C53030}
\definecolor{riskamber}{HTML}{D69E2E}
\definecolor{riskgreen}{HTML}{38A169}
\definecolor{nvgreen}{HTML}{76B900}

% ============ 表格 ============
\usepackage{booktabs}
\usepackage{longtable}
\usepackage{tabularx}
\usepackage{multirow}
\usepackage{array}
\newcolumntype{Y}{>{\centering\arraybackslash}X}
\newcolumntype{L}[1]{>{\raggedright\arraybackslash}p{#1}}
\newcolumntype{C}[1]{>{\centering\arraybackslash}p{#1}}
\newcolumntype{R}[1]{>{\raggedleft\arraybackslash}p{#1}}

% ============ 图形 ============
\usepackage{tikz}
\usetikzlibrary{shapes,arrows.meta,positioning,calc,backgrounds,fit}
\usepackage{pgfplots}
\usepgfplotslibrary{polar}
\pgfplotsset{compat=1.18}
\usepackage[edges]{forest}

% ============ 页眉页脚（报告标题和日期需在main.tex中覆盖） ============
\usepackage{fancyhdr}
\pagestyle{fancy}
\fancyhf{}
\fancyhead[L]{\small\color{steelgrey} \reporttitle}
\fancyhead[R]{\small\color{steelgrey} \reportdate}
\fancyfoot[C]{\small\color{steelgrey} \thepage}
\fancyfoot[R]{\tiny\color{steelgrey} 本报告不构成任何投资建议}
\renewcommand{\headrulewidth}{0.4pt}
\renewcommand{\footrulewidth}{0.2pt}

% ============ 章节格式 ============
\usepackage{titlesec}
\titleformat{\chapter}[display]
  {\normalfont\huge\bfseries\color{navy}}
  {\chaptertitlename\ \thechapter}{20pt}{\Huge}
\titleformat{\section}
  {\normalfont\Large\bfseries\color{navy}}
  {\thesection}{1em}{}
\titleformat{\subsection}
  {\normalfont\large\bfseries\color{steelgrey}}
  {\thesubsection}{1em}{}

% ============ 超链接 ============
\usepackage{hyperref}
\hypersetup{
  colorlinks=true,
  linkcolor=navy,
  urlcolor=navy,
  citecolor=navy,
}

% ============ 其他工具 ============
\usepackage{enumitem}
\setlist{nosep, leftmargin=2em}
\usepackage{fontawesome5}
\usepackage{amssymb}
\usepackage{pifont}
\usepackage{tcolorbox}
\tcbuselibrary{skins,breakable}
\usepackage{caption}
\captionsetup{font=small, skip=4pt}
\usepackage{float}
\setlength{\textfloatsep}{8pt plus 2pt minus 2pt}
\setlength{\floatsep}{6pt plus 2pt minus 2pt}
\setlength{\intextsep}{6pt plus 2pt minus 2pt}
\setlength{\abovecaptionskip}{4pt}
\setlength{\belowcaptionskip}{2pt}

% ============ 自定义命令 ============
\newcommand{\rmark}{\textcolor{riskred}{\ding{108}}}
\newcommand{\amark}{\textcolor{riskamber}{\ding{108}}}
\newcommand{\gmark}{\textcolor{riskgreen}{\ding{108}}}
\newcommand{\starfive}{$\bigstar\bigstar\bigstar\bigstar\bigstar$}
\newcommand{\starfour}{$\bigstar\bigstar\bigstar\bigstar$}
\newcommand{\starthree}{$\bigstar\bigstar\bigstar$}

% 信息框
\newtcolorbox{keyinsight}[1][]{
  colback=lightgrey, colframe=navy,
  fonttitle=\bfseries, title={#1}, breakable
}
\newtcolorbox{riskbox}[1][]{
  colback=red!5, colframe=riskred,
  fonttitle=\bfseries, title={#1}, breakable
}
 after \documentclass[a4paper,11pt,openany,fontset=none]{ctexrep}

% ============ 页面布局 ============
\usepackage[top=2cm, bottom=2cm, left=2cm, right=2cm]{geometry}
\setlength{\parskip}{0.3em}
\setlength{\parindent}{2em}
\renewcommand{\baselinestretch}{1.15}

% ============ 字体（macOS） ============
\setCJKmainfont[BoldFont={Heiti SC}]{STSong}
\setCJKsansfont{Heiti SC}
\setCJKmonofont{Heiti SC}
\setmainfont{Times New Roman}

% ============ 颜色（IB Navy/Grey palette） ============
\usepackage{xcolor}
\definecolor{navy}{HTML}{003366}
\definecolor{steelgrey}{HTML}{4A5568}
\definecolor{lightgrey}{HTML}{F7FAFC}
\definecolor{riskred}{HTML}{C53030}
\definecolor{riskamber}{HTML}{D69E2E}
\definecolor{riskgreen}{HTML}{38A169}
\definecolor{nvgreen}{HTML}{76B900}

% ============ 表格 ============
\usepackage{booktabs}
\usepackage{longtable}
\usepackage{tabularx}
\usepackage{multirow}
\usepackage{array}
\newcolumntype{Y}{>{\centering\arraybackslash}X}
\newcolumntype{L}[1]{>{\raggedright\arraybackslash}p{#1}}
\newcolumntype{C}[1]{>{\centering\arraybackslash}p{#1}}
\newcolumntype{R}[1]{>{\raggedleft\arraybackslash}p{#1}}

% ============ 图形 ============
\usepackage{tikz}
\usetikzlibrary{shapes,arrows.meta,positioning,calc,backgrounds,fit}
\usepackage{pgfplots}
\usepgfplotslibrary{polar}
\pgfplotsset{compat=1.18}
\usepackage[edges]{forest}

% ============ 页眉页脚（报告标题和日期需在main.tex中覆盖） ============
\usepackage{fancyhdr}
\pagestyle{fancy}
\fancyhf{}
\fancyhead[L]{\small\color{steelgrey} \reporttitle}
\fancyhead[R]{\small\color{steelgrey} \reportdate}
\fancyfoot[C]{\small\color{steelgrey} \thepage}
\fancyfoot[R]{\tiny\color{steelgrey} 本报告不构成任何投资建议}
\renewcommand{\headrulewidth}{0.4pt}
\renewcommand{\footrulewidth}{0.2pt}

% ============ 章节格式 ============
\usepackage{titlesec}
\titleformat{\chapter}[display]
  {\normalfont\huge\bfseries\color{navy}}
  {\chaptertitlename\ \thechapter}{20pt}{\Huge}
\titleformat{\section}
  {\normalfont\Large\bfseries\color{navy}}
  {\thesection}{1em}{}
\titleformat{\subsection}
  {\normalfont\large\bfseries\color{steelgrey}}
  {\thesubsection}{1em}{}

% ============ 超链接 ============
\usepackage{hyperref}
\hypersetup{
  colorlinks=true,
  linkcolor=navy,
  urlcolor=navy,
  citecolor=navy,
}

% ============ 其他工具 ============
\usepackage{enumitem}
\setlist{nosep, leftmargin=2em}
\usepackage{fontawesome5}
\usepackage{amssymb}
\usepackage{pifont}
\usepackage{tcolorbox}
\tcbuselibrary{skins,breakable}
\usepackage{caption}
\captionsetup{font=small, skip=4pt}
\usepackage{float}
\setlength{\textfloatsep}{8pt plus 2pt minus 2pt}
\setlength{\floatsep}{6pt plus 2pt minus 2pt}
\setlength{\intextsep}{6pt plus 2pt minus 2pt}
\setlength{\abovecaptionskip}{4pt}
\setlength{\belowcaptionskip}{2pt}

% ============ 自定义命令 ============
\newcommand{\rmark}{\textcolor{riskred}{\ding{108}}}
\newcommand{\amark}{\textcolor{riskamber}{\ding{108}}}
\newcommand{\gmark}{\textcolor{riskgreen}{\ding{108}}}
\newcommand{\starfive}{$\bigstar\bigstar\bigstar\bigstar\bigstar$}
\newcommand{\starfour}{$\bigstar\bigstar\bigstar\bigstar$}
\newcommand{\starthree}{$\bigstar\bigstar\bigstar$}

% 信息框
\newtcolorbox{keyinsight}[1][]{
  colback=lightgrey, colframe=navy,
  fonttitle=\bfseries, title={#1}, breakable
}
\newtcolorbox{riskbox}[1][]{
  colback=red!5, colframe=riskred,
  fonttitle=\bfseries, title={#1}, breakable
}
 after \documentclass[a4paper,11pt,openany,fontset=none]{ctexrep}

% ============ 页面布局 ============
\usepackage[top=2cm, bottom=2cm, left=2cm, right=2cm]{geometry}
\setlength{\parskip}{0.3em}
\setlength{\parindent}{2em}
\renewcommand{\baselinestretch}{1.15}

% ============ 字体（macOS） ============
\setCJKmainfont[BoldFont={Heiti SC}]{STSong}
\setCJKsansfont{Heiti SC}
\setCJKmonofont{Heiti SC}
\setmainfont{Times New Roman}

% ============ 颜色（IB Navy/Grey palette） ============
\usepackage{xcolor}
\definecolor{navy}{HTML}{003366}
\definecolor{steelgrey}{HTML}{4A5568}
\definecolor{lightgrey}{HTML}{F7FAFC}
\definecolor{riskred}{HTML}{C53030}
\definecolor{riskamber}{HTML}{D69E2E}
\definecolor{riskgreen}{HTML}{38A169}
\definecolor{nvgreen}{HTML}{76B900}

% ============ 表格 ============
\usepackage{booktabs}
\usepackage{longtable}
\usepackage{tabularx}
\usepackage{multirow}
\usepackage{array}
\newcolumntype{Y}{>{\centering\arraybackslash}X}
\newcolumntype{L}[1]{>{\raggedright\arraybackslash}p{#1}}
\newcolumntype{C}[1]{>{\centering\arraybackslash}p{#1}}
\newcolumntype{R}[1]{>{\raggedleft\arraybackslash}p{#1}}

% ============ 图形 ============
\usepackage{tikz}
\usetikzlibrary{shapes,arrows.meta,positioning,calc,backgrounds,fit}
\usepackage{pgfplots}
\usepgfplotslibrary{polar}
\pgfplotsset{compat=1.18}
\usepackage[edges]{forest}

% ============ 页眉页脚（报告标题和日期需在main.tex中覆盖） ============
\usepackage{fancyhdr}
\pagestyle{fancy}
\fancyhf{}
\fancyhead[L]{\small\color{steelgrey} \reporttitle}
\fancyhead[R]{\small\color{steelgrey} \reportdate}
\fancyfoot[C]{\small\color{steelgrey} \thepage}
\fancyfoot[R]{\tiny\color{steelgrey} 本报告不构成任何投资建议}
\renewcommand{\headrulewidth}{0.4pt}
\renewcommand{\footrulewidth}{0.2pt}

% ============ 章节格式 ============
\usepackage{titlesec}
\titleformat{\chapter}[display]
  {\normalfont\huge\bfseries\color{navy}}
  {\chaptertitlename\ \thechapter}{20pt}{\Huge}
\titleformat{\section}
  {\normalfont\Large\bfseries\color{navy}}
  {\thesection}{1em}{}
\titleformat{\subsection}
  {\normalfont\large\bfseries\color{steelgrey}}
  {\thesubsection}{1em}{}

% ============ 超链接 ============
\usepackage{hyperref}
\hypersetup{
  colorlinks=true,
  linkcolor=navy,
  urlcolor=navy,
  citecolor=navy,
}

% ============ 其他工具 ============
\usepackage{enumitem}
\setlist{nosep, leftmargin=2em}
\usepackage{fontawesome5}
\usepackage{amssymb}
\usepackage{pifont}
\usepackage{tcolorbox}
\tcbuselibrary{skins,breakable}
\usepackage{caption}
\captionsetup{font=small, skip=4pt}
\usepackage{float}
\setlength{\textfloatsep}{8pt plus 2pt minus 2pt}
\setlength{\floatsep}{6pt plus 2pt minus 2pt}
\setlength{\intextsep}{6pt plus 2pt minus 2pt}
\setlength{\abovecaptionskip}{4pt}
\setlength{\belowcaptionskip}{2pt}

% ============ 自定义命令 ============
\newcommand{\rmark}{\textcolor{riskred}{\ding{108}}}
\newcommand{\amark}{\textcolor{riskamber}{\ding{108}}}
\newcommand{\gmark}{\textcolor{riskgreen}{\ding{108}}}
\newcommand{\starfive}{$\bigstar\bigstar\bigstar\bigstar\bigstar$}
\newcommand{\starfour}{$\bigstar\bigstar\bigstar\bigstar$}
\newcommand{\starthree}{$\bigstar\bigstar\bigstar$}

% 信息框
\newtcolorbox{keyinsight}[1][]{
  colback=lightgrey, colframe=navy,
  fonttitle=\bfseries, title={#1}, breakable
}
\newtcolorbox{riskbox}[1][]{
  colback=red!5, colframe=riskred,
  fonttitle=\bfseries, title={#1}, breakable
}


\begin{document}

\begin{center}
{\sffamily\small\bfseries\color{steelgrey} \reportkicker}\\[3pt]
{\LARGE\bfseries\color{navy} \reporttitle}\\[5pt]
{\large\color{steelgrey} \reportsubtitle}\\[4pt]
{\small\color{mutedgrey} \reportscope \quad|\quad \reportdate \quad|\quad \reportdatacutoff}
\end{center}
\vspace{0.7em}

\begin{dashboardbox}[Decision Snapshot]
\begin{houseviewbox}[House View]
\reporthouseview
\end{houseviewbox}
\begin{sourcequalitybox}[Data Quality]
\reportquality
\end{sourcequalitybox}
\end{dashboardbox}

%% ================================================================
\section{市场格局：高低切换已成定局}

\begin{exhibitbox}[Exhibit 1: 7月初资金迁徙全景]
\begin{tabularx}{\textwidth}{L{4cm}X}
\toprule
\textbf{信号} & \textbf{数据} \\
\midrule
北向资金 & 7月整体小幅净流出90亿，但\textbf{内部极端分化}：抛售高位算力中游兑现80亿+，同步加仓半导体设备/存储/电子特气/封装合计净流入110亿+ \\
7/7北向 & 当日净流入98.44亿，重点低位硬科技+高股息 \\
龙虎榜 & 机构/社保/QFII持续扫货半导体设备+创新药低位龙头 \\
高位兑现 & AI算力、光模块题材股机构席位净卖出12亿+（7/7当日） \\
板块资金7/6 & 计算机净流入第一；医药、农牧获流入；电子/通信净流出 \\
迁徙周期 & 历史规律：7月上旬机构系统性调仓持续15--20个交易日 \\
\bottomrule
\end{tabularx}
\sourcenote{东方财富、同花顺龙虎榜、凤凰网财经（2026.07.06-07）}
\end{exhibitbox}

\textbf{核心判断}：这不是一日游行情，而是机构系统性的高低切换，预计持续至7月下旬。低位板块正处于主升浪启动的早期阶段。

%% ================================================================
\section{板块一：创新药（已启动主升，最确定）}

\subsection*{核心信号}
\begin{itemize}
  \item 板块指数6/29触底1101点后连续反弹，走出独立主升行情。
  \item \textbf{主力连续3日净流入超10亿}（截至7/6），板块单日净流入32.99亿。
  \item 估值处近5年\textbf{20\%分位以下}，历史低位区域。
  \item 北向+公募+私募+产业资本6月底起持续大手笔加仓。
  \item 多条均线站稳，KDJ从超卖区金叉，放量上涨确认趋势。
  \item 6/29至7/6区间：热景生物累涨60\%+，海南海药/苑东生物/石药创新等6只涨超30\%。
\end{itemize}

\subsection*{驱动逻辑}
\begin{enumerate}
  \item 两年持续调整后套牢盘消化充分，上方抛压极轻。
  \item 政策面利好密集（医保谈判规则优化、创新药审批加速）。
  \item 中报业绩窗口：多家创新药企营收/管线里程碑兑现。
  \item 资金面：上半年公募严重低配医药，存在回补需求。
\end{enumerate}

\subsection*{重点标的}
\begin{tabularx}{\textwidth}{L{3cm}X}
\toprule
\textbf{个股} & \textbf{逻辑} \\
\midrule
甘李药业 & 创新药+胰岛素龙头，低位主力持续吸筹 \\
恒瑞医药 & 创新药一哥，北向重仓，估值修复空间大 \\
百济神州 & 全球化创新药龙头，管线里程碑密集 \\
石药创新 & 区间涨幅30\%+，仍处低位，肿瘤管线丰富 \\
苑东生物 & 区间涨幅30\%+，中报预增 \\
\bottomrule
\end{tabularx}

%% ================================================================
\section{板块二：半导体设备/材料（北向110亿吸筹，长线核心）}

\subsection*{核心信号}
\begin{itemize}
  \item \textbf{北向7月逆势加仓半导体上游合计114亿元}，单只净流入4--19亿。
  \item 海外主权基金连续半年稳步加仓锁仓。
  \item 高盛7月初发布深度研报唱多国内半导体设备。
  \item 大基金三期持续注资设备材料端。
  \item 头部设备企业在手订单全部顺延至2027年，Q1净利同比增幅近200\%。
  \item 全年横盘低位，筹码充分集中，具备突破条件。
\end{itemize}

\subsection*{驱动逻辑}
\begin{enumerate}
  \item 国内晶圆厂（中芯/华虹/长存）扩产订单锁定至2027，业绩增长确定性极强。
  \item 国产化率持续提升（设备国产化从20\%向40\%+推进）。
  \item AI芯片产能扩张→设备/材料为最确定受益环节。
  \item Chiplet/2.5D/3D混合键合→先进封装设备全新增量。
  \item 区别于下游存储/算力已翻倍，上游设备全年横盘，估值洼地。
\end{enumerate}

\subsection*{重点标的}
\begin{tabularx}{\textwidth}{L{3cm}X}
\toprule
\textbf{个股} & \textbf{逻辑} \\
\midrule
北方华创 & 设备龙头，北向持续加仓4--19亿，大基金重仓 \\
中微公司 & 刻蚀设备龙头，北向加仓，订单饱满 \\
拓荆科技 & 薄膜沉积设备，订单延至2027 \\
富创精密 & 设备零部件，北向加仓 \\
安集科技 & CMP材料龙头，北向持续流入 \\
雅克科技 & 电子特气/光刻胶，北向加仓 \\
芯碁微装 & CoPoS板级光刻龙头，国产首台量产 \\
\bottomrule
\end{tabularx}

%% ================================================================
\section{板块三：军工（估值10年最低，社保/北向逆势建仓）}

\subsection*{核心信号}
\begin{itemize}
  \item 整体估值处于\textbf{近10年18\%低位区间}，A股少有的``低估值+高增长''。
  \item 军工ETF连续3日吸金1.09亿+，规模达69.59亿。
  \item \textbf{社保二季度大幅加仓}军工上游材料、航空零部件龙头。
  \item 北向资金持续回流，锁定12只核心标的。
  \item 公募军工配置比例处近5年低位，后续有充足加仓空间。
\end{itemize}

\subsection*{驱动逻辑}
\begin{enumerate}
  \item 十五五国防建设规划稳步推进，军费投入稳定增长。
  \item 军贸出口爆发（无人机、高端装备出口订单大幅增长）。
  \item 订单饱满、产能充足，全年业绩增长确定性高。
  \item 前期市场关注度极低，估值压制充分，弹簧压到极致。
  \item 7/6新规落地催化：长线资金加速完成加仓。
\end{enumerate}

\subsection*{重点标的}
\begin{tabularx}{\textwidth}{L{3cm}X}
\toprule
\textbf{个股} & \textbf{逻辑} \\
\midrule
航发动力 & 航空发动机核心，社保/北向重仓 \\
中航沈飞 & 战斗机整机，订单确定性 \\
中航光电 & 军工连接器龙头，北向加仓 \\
抚顺特钢 & 高温合金材料，社保加仓 \\
紫光国微 & 军工电子/FPGA，底部筹码集中 \\
\bottomrule
\end{tabularx}

%% ================================================================
\section{板块四：人形机器人零部件（超跌批量涨停，高弹性）}

\subsection*{核心信号}
\begin{itemize}
  \item 7/3板块爆发：超40只涨停，板块净流入170亿+，成交5300亿。
  \item 7月初机器人概念单日主力净流入143.74亿，人形机器人细分净流入113.78亿（全市场第一）。
  \item 多数标的为前期超跌小票（流通市值二三十亿），少量资金即可推动涨停潮。
  \item 产业催化密集：宇树科技IPO + 特斯拉Optimus量产消息。
\end{itemize}

\subsection*{驱动逻辑}
\begin{enumerate}
  \item 人形机器人产业从0到1，2026年为量产元年。
  \item 减速器、伺服电机、力矩执行器为核心运动部件，技术壁垒高。
  \item 产业政策持续落地（国家级扶持）。
  \item 超跌后弹性最大，适合短线博弈+中线布局。
  \item 三花智控7/6尾盘被抢筹21.6亿（板块中最大单笔）。
\end{enumerate}

\subsection*{重点标的}
\begin{tabularx}{\textwidth}{L{3cm}X}
\toprule
\textbf{个股} & \textbf{逻辑} \\
\midrule
三花智控 & 热管理+机器人温控双龙头，7/6抢筹21.6亿 \\
雷赛智能 & 伺服执行器龙头，4连板 \\
绿的谐波 & 谐波减速器龙头 \\
万里扬 & 减速器，涨停 \\
科瑞技术 & 机器人精密零部件，股东户数降17.65\%（主力重度吸筹） \\
凯龙高科 & 执行器，20cm涨停 \\
\bottomrule
\end{tabularx}

%% ================================================================
\section{板块五：氟化工/电子特气（涨价+AI绑定，周期反转）}

\subsection*{核心信号}
\begin{itemize}
  \item 产品价格全线暴涨：R32从1.5万$\to$6.4万/吨（+327\%），六氟化钨+232\%，电子氢氟酸+17--19\%。
  \item 北向持续加仓低位电子特气标的。
  \item 行业定性从``周期化工''转型为``AI+半导体战略新材料''。
  \item 国家配额制约束供给（萤石不能多挖），底价被彻底抬高。
\end{itemize}

\subsection*{驱动逻辑}
\begin{enumerate}
  \item AI芯片产能爆发→晶圆制造用电子特气需求暴增（特气占芯片制造成本13--16\%）。
  \item 海外高端材料供给收缩+国产替代加速。
  \item 三代制冷剂配额政策刚性（2025年起配额锁死，供给不会新增）。
  \item 下游对价格不敏感（芯片厂不会因特气涨价停产）。
  \item 瑞银研报定性：``AI为氟化工创造万亿级全新增量市场''。
\end{enumerate}

\subsection*{重点标的}
\begin{tabularx}{\textwidth}{L{3cm}X}
\toprule
\textbf{个股} & \textbf{逻辑} \\
\midrule
巨化股份 & 氟化工龙头，三代配额优势 \\
三美股份 & 制冷剂+电子特气，配额龙头 \\
昊华科技 & 高端氟材料+电子气体 \\
雅克科技 & 电子特气（半导体级），北向加仓 \\
金石资源 & 萤石（氟化工上游），配额稀缺 \\
\bottomrule
\end{tabularx}

%% ================================================================
\section{筹码集中度异动（主力悄悄吸筹信号）}

截至6月30日，101只A股股东户数连续3期以上下降（最多连降14期），呈现持续筹码集中。核心特征：低位横盘+股东户数降>10\%+成交量温和=机构分批承接。

\textbf{近期筹码集中典型个股}：
\begin{itemize}
  \item \textbf{科瑞技术}（002957）：股东户数降17.65\%（最大降幅之一），机器人精密零部件，低位横盘后主力重度吸筹。
  \item \textbf{鼎龙科技}：股东户数降3.37\%，半导体CMP材料。
  \item 101只筹码连续集中股中，半导体设备/材料、创新药、军工电子占比最高。
\end{itemize}

%% ================================================================
\section{综合排序与操作建议}

\begin{tabularx}{\textwidth}{L{0.5cm}L{3cm}L{2cm}L{2cm}L{2.5cm}X}
\toprule
 & \textbf{板块} & \textbf{启动阶段} & \textbf{确定性} & \textbf{弹性} & \textbf{适合风格} \\
\midrule
1 & 创新药 & \textbf{已启动主升} & 极高 & 中高 & 趋势跟随，可追 \\
2 & 半导体设备 & 底部蓄势末期 & 极高 & 中 & 长线底仓，等突破 \\
3 & 军工 & 底部蓄势中 & 高 & 中高 & 左侧布局，耐心持有 \\
4 & 人形机器人 & 脉冲启动（高波动） & 中 & 极高 & 短线博弈+中线核心 \\
5 & 氟化工/电子特气 & 已启动（涨价确认） & 高 & 高 & 周期+成长，可跟 \\
\bottomrule
\end{tabularx}

\subsection*{操作建议}

\begin{itemize}
  \item \textbf{稳健型}：创新药（已确认主升）+ 半导体设备（北向锁仓、等突破右侧确认）。
  \item \textbf{进攻型}：人形机器人零部件（弹性大但波动大）+ 氟化工（涨价驱动业绩弹性）。
  \item \textbf{长线配置}：军工（估值极低、订单确定、社保在建仓）+ 半导体设备。
  \item \textbf{回避}：上半年翻倍的AI算力/光模块/高位题材股（资金正在系统性流出）。
\end{itemize}

\subsection*{时间窗口}

高位科技兑现资金的迁徙周期历史规律为\textbf{15--20个交易日}（7月上旬启动$\to$7月下旬完成）。当前处于迁徙早期（第3--5天），低位板块仍有充足上行空间。但7月下旬迁徙完成后可能出现板块内部分化，届时需聚焦中报业绩确认的个股。

%% ================================================================
\vfill
\begin{disclosurebox}[Metadata]
\footnotesize
Confidence: 88/100 \quad | \quad
Data quality: high / multi-source cross-validated \quad | \quad
Degradation: None material.

\vspace{0.3em}
Remaining uncertainty (12 pts): (1) 迁徙周期是否因外部事件缩短/延长不确定；(2) 个股层面需中报业绩最终确认；(3) 地缘/政策黑天鹅不可预测。

\vspace{0.4em}
{\tiny\reportdisclaimer\ Generated by AI agents for research reference only. 2026-07-07.}
\end{disclosurebox}

\end{document}
